\section{Introduction: replacing execution, preserving the function}

Before an enterprise buys software, connects systems, or releases an application, people review
its architecture, security, contracts, cost, and operational readiness. The required outputs are
decisions and their supporting work: identify a missing agreement, reject an unsafe network
design, require a recovery test, or authorize a conditional release. These functions persist
even if software performs work previously assigned to analysts, architects, and committee
support staff. A professional title is an organizational allocation of tasks, not a technical
specification of how those tasks must be executed.

This paper argues that a Digital Governance Framework (DGF), understood as an enterprise's
connected review gates, is a tractable candidate for early substitution of human review tasks
by agentic systems. The claim concerns execution: an agent that produces a draft for someone
to redo has supplied assistance; a system that completes an accepted review has performed that
task. A workflow can contain substituted tasks and remaining human decisions at the same time.
The relevant system can combine language models, deterministic software, and escalation.
Replacing every rule check with a language-model call is neither necessary nor desirable.

Figure~\ref{fig:dgf-routes} makes this organization concrete: the same project dossier moves
through several specialist reviews, each producing a decision and its supporting record.
Buying a solution, connecting systems, and building an application require different review
sequences. General consolidates the preceding reviews into the route's governance record.

\begin{figure}[!htb]
\centering
\begin{tikzpicture}[
  x=2.5cm,y=1cm,
  gate/.style={draw=black!55,rounded corners=2pt,minimum height=0.64cm,
    text width=2.08cm,align=center,inner sep=2pt,font=\small},
  consolidate/.style={gate,fill=black!8,draw=black!75},
  flow/.style={-{Stealth[length=1.5mm]},draw=black!65,semithick},
  route/.style={anchor=west,font=\small}
]
\node[route] at (-0.46,0.62) {\textbf{BUY} \quad Purchase a solution};
\node[gate,fill=blue!5] (b1) at (0,0) {Procurement};
\node[gate,fill=blue!5] (b2) at (1,0) {Legal};
\node[gate,fill=blue!5] (b3) at (2,0) {Compliance};
\node[gate,fill=blue!5] (b4) at (3,0) {Security};
\node[gate,fill=blue!5] (b5) at (4,0) {IT};
\node[consolidate] (b6) at (5,0) {General};
\foreach \a/\b in {b1/b2,b2/b3,b3/b4,b4/b5,b5/b6}
  \draw[flow] (\a) -- (\b);

\node[route] at (-0.46,-0.73) {\textbf{INTEGRATE} \quad Connect existing systems};
\node[gate,fill=teal!6] (i1) at (0,-1.35) {IT};
\node[gate,fill=teal!6] (i2) at (1,-1.35) {Architecture};
\node[gate,fill=teal!6] (i3) at (2,-1.35) {Security};
\node[gate,fill=teal!6] (i4) at (3,-1.35) {Legal};
\node[gate,fill=teal!6] (i5) at (4,-1.35) {Compliance};
\node[consolidate] (i6) at (5,-1.35) {General};
\foreach \a/\b in {i1/i2,i2/i3,i3/i4,i4/i5,i5/i6}
  \draw[flow] (\a) -- (\b);

\node[route] at (-0.46,-2.08) {\textbf{BUILD} \quad Develop an application};
\node[gate,fill=orange!7] (d1) at (0,-2.7) {IT};
\node[gate,fill=orange!7] (d2) at (1.25,-2.7) {Architecture};
\node[gate,fill=orange!7] (d3) at (2.5,-2.7) {Security};
\node[gate,fill=orange!7] (d4) at (3.75,-2.7) {Tech\\Readiness};
\node[consolidate] (d5) at (5,-2.7) {General};
\foreach \a/\b in {d1/d2,d2/d3,d3/d4,d4/d5}
  \draw[flow] (\a) -- (\b);

\node[font=\small,align=center] at (2.5,-3.48)
  {\textbf{At each gate:} inspect the dossier; record findings, evidence, actions, and a decision.};
\end{tikzpicture}
\caption{Three DGF workflows evaluated in DGF-Bench. A project follows the route matching its
need. Arrows show review order and information handoffs; General consolidates preceding reviews.
The benchmark runs all scheduled gates, including after a refusal. These are example governance
routes, not automatic approvals or a universal enterprise standard.}
\label{fig:dgf-routes}
\end{figure}


The DGF is attractive because it already organizes work into dossiers, policies, named
authorities, bounded decisions, and recorded handoffs. This gives engineering a starting point
that open-ended management lacks. The same features also reveal why a fluent agent is insufficient:
facts may be missing, exceptions expensive, evidence unfaithful, and commitments unauthorized.
Our claim is that these are explicit conditions of substitution, rather than reasons to assume
that the current human task allocation must remain permanent. Whether DGF is earlier or easier
to automate than every other business function is a comparative hypothesis, not an observed ranking.

We connect three contributions. First, a gate-contract formulation identifies information,
validity, and authority requirements for delegating review execution. A source counterexample
shows when no reviewer can satisfy a fixed decision contract from the permitted information.
Second, a complete labor account gives a substitution threshold: reducing routine execution
only reduces required human work if exception handling, verification, correction, and support
remain below a computable budget. Third, DGF-Bench provides a public, controlled measurement
of specified review tasks and their failure modes, including a deterministic control and
repeated trajectories. The purpose of linking them is to make the route from an agent score
to a workforce claim inspectable; none of the three substitutes for the others.

\subsection{Research questions and units of analysis}

We organize the argument around three questions with different observational units.
\emph{First, can software execute a specified gate contract?} Its unit is a review occurrence,
with fixed inputs, rules, and acceptance conditions. DGF-Bench measures this question in a
synthetic environment. \emph{Second, can that execution remain dependable across a project?}
Its unit is the complete route, including the evidence and conditions passed between gates.
Route success and repeated trajectories address parts of this question. \emph{Third, does
deployment reduce the human work required to deliver the governed output?} Its unit is an
operating population over a stated period. Answering it requires a labor account, including
work added outside the nominal review team. This third quantity is not present in an API trace.

The units cannot be exchanged without assumptions. Five thousand scored gates are not five
thousand independent companies. A successful run on a generated dossier is not a measured
number of hours saved. A refusal may be a successful review even though the underlying project
does not proceed. Defining these distinctions makes a positive replacement claim more precise:
software can execute an accepted task while humans remain responsible for setting its scope,
maintaining its infrastructure, or resolving a subset of exceptions.

\begin{table}[!htb]
\centering\small
\caption{Claim structure and evidence required. A result at one level does not supply the
missing observation at another.}
\label{tab:claim-structure}
\begin{tabularx}{\linewidth}{@{}P{2.8cm}LL@{}}
\toprule
Claim & Observable test & Evidence in this paper\\
\midrule
Contract execution & Decision, findings, actions, evidence, and mandate meet fixed criteria & Original runs and deterministic control\\
Workflow reliability & Every required gate succeeds; repeated trajectories remain successful & Complete-route scores and 135 additional runs\\
Net task substitution & Accepted output requires fewer total human hours & Accounting conditions and illustrative calculations\\
Ecosystem workforce reduction & Comparable governed output and quality with a lower complete FTE account & Dated hypothesis and a measurement protocol\\
\bottomrule
\end{tabularx}
\end{table}

The central thesis is consequently about the implementation of a function. Organizations may
still require architecture assurance, a security judgment, or evidence of operational readiness.
That continuing requirement does not establish that every dossier needs the same manual
execution. At the same time, describing a function as a contract does not make its evidence
available or resolve policy disagreement. The proposed engineering work is to turn the subset
with adequate information and delegated authority into executable review services, and then
measure the cost of the boundary that remains.

The paper proceeds from the contract to the labor account, then to its FDE implementation and
experimental evidence. We conclude with a deployment test for the workforce hypothesis.
This manuscript synthesizes the longer \emph{The Last Human Gate} and replaces the previous
empirical companion in the repository. The versions share their experimental record;
the expanded exposition reports no additional model calls or enterprise observations.
